% titel: Blatt 0, A9: Potenzen

% Teil I: Werkzeuge, kein Lösungsweg.

Ausdrücke mit Potenzen und Wurzeln ohne Taschenrechner bestimmen bzw.
vereinfachen.

\begin{werkzeugkasten}[Potenzgesetze]
Für $a>0$ und rationale Exponenten:
\begin{align}
  a^{-n}&=\frac1{a^n}, &
  a^{\frac mn}&=\sqrt[n]{a^m}, \quad a>0, \\
  a^r a^s&=a^{r+s}, &
  \frac{a^r}{a^s}&=a^{r-s}, &
  (a^r)^s&=a^{rs}, &
  \Bigl(\frac ab\Bigr)^r&=\frac{a^r}{b^r} .
\end{align}
Eine Wurzel ist eine Potenz mit gebrochenem Exponenten;
$\sqrt a=a^{1/2}$, $\sqrt[n]a=a^{1/n}$.
\end{werkzeugkasten}

\begin{werkzeugkasten}[Gebrochene und negative Exponenten]
Umkehren und Wurzel ziehen statt ``im Kopf'' rechnen:
\begin{equation}
  27^{-1/3}=\frac1{27^{1/3}}=\frac13,\qquad
  \Bigl(\frac49\Bigr)^{-1/2}=\Bigl(\frac94\Bigr)^{1/2}=\frac32 .
\end{equation}
Negative Exponenten wechseln in den Nenner; eine Potenz mit
`${}^{-1/2}$` kehrt den Bruch um und zieht die Wurzel.
\end{werkzeugkasten}

\begin{achtung}[Darauf achten]
\begin{itemize}
  \item $a^{-n}\neq -a^n$. Der Minus-Exponent bedeutet Kehrwert.
  \item Beim Vereinfachen von $a^5 a^{-2}/a^3$ die Exponenten addieren bzw. subtrahieren.
  \item Bruch-Hochzahlen: erst Wurzel (oder umgekehrt), z. B.
        $\bigl(x^{2/3}\bigr)^{-3}=x^{-2}=\tfrac1{x^2}$ -- das Ergebnis ohne
        negativen Exponenten schreiben.
  \item $a^{3/4}$ als Wurzel: $a^{3/4}=\sqrt[4]{a^3}$.
\end{itemize}
\end{achtung}
